\section{Implementation}\label{sec:impl}

\subsection{Extended fragment and front-end}
\label{sec:impl-fragment}
\textsc{TensorGuard}'s analyser parses class source via Python's
\texttt{ast} module, extracts a typed IR, and dispatches each
op to a handler that returns the refinement-typed result type.
We add four front-end families under \texttt{src/v5/} (no edits
to legacy \texttt{src/}). The \emph{symbolic config} environment
binds \texttt{self.config.<attr>} to fresh \texttt{SymInt}s so
that HuggingFace-style modules whose dimensions come from a
\texttt{config} object can be checked once and reused. The
\emph{QKV unpacking} front-end recognises the four idiomatic
Q/K/V splits (\texttt{split}, \texttt{chunk},
\texttt{view+unbind}, and \texttt{einops.rearrange} with the
canonical \texttt{(three h d) $\to$ three b h t d} pattern) and
elaborates each into three correlated tensor types. The
\emph{negative-one reshape} elaborator emits a Z3 divisibility
constraint $P\bmod Q=0$. Transfer rules for
\texttt{F.scaled\_dot\_product\_attention},
\texttt{nn.MultiheadAttention}, \texttt{nn.LayerNorm}, and
\texttt{nn.RMSNorm} register into the analyser's shape-op
dictionary. All $74$ unit tests pass on \texttt{python3.11} +
\texttt{torch~2.9.1}. We report honestly that on the \emph{legacy
v4 corpus} these new APIs produce zero verdict-count delta (the
v4 corpus did not abstain on patterns these APIs target); the
v5 contribution is that the new APIs are now \emph{callable on
realistic Transformer modules} that v4 could not enter, surfaced
in the empirical evaluation through the 488-block real-source
corpus (\Cref{sec:eval-benchmark}).

\subsection{Backward verifier}\label{sec:impl-backward}
The backward verifier (\texttt{src/v5/backward\_shape.py},
\texttt{grad\_flag\_verifier.py}) consumes a forward graph
annotated with the refinement types of \Cref{sec:calculus} and
certifies, without execution, that (i)~\texttt{loss.backward()}
would produce gradients of the right shape modulo
\texttt{sum\_to(grad,primal.shape)} reduction, and (ii)~the set
$H$ of parameters reverse-reachable from \texttt{loss} along
non-detached, non-\texttt{no\_grad}-context,
\texttt{requires\_grad=True} edges equals the user-declared
expected-to-learn set $E$. The check is a reverse-topological
traversal mirroring PyTorch's per-Function shape contract
\verb|grad_outputs[i].shape == y_i.shape|. The grad-flag
component additionally detects the three canonical
silent-zero-grad bug classes: in-place writes on a leaf with
\texttt{requires\_grad=True}; no leaf in the trace has
\texttt{requires\_grad=True}; and parameters used only inside
\texttt{torch.no\_grad()} or after \texttt{.detach()}. These
are reported as \textsc{Refuted} with witness traces.

\paragraph{Empirical agreement.}
On a $500$-model property test (random small \texttt{nn.Module}s
drawn from a grammar) the static verifier agrees with PyTorch's
runtime on $500/500$ models (gradient presence and gradient
shape both match). On $8$ hand-curated case studies modelled on
canonical PyTorch issues (\texttt{no\_grad} scope,
\texttt{.detach()}, \texttt{requires\_grad=False},
in-place-on-leaf, missing leaf, in-place alias mutating a saved
tensor, non-scalar loss, unused parameter silently skipped by
\texttt{optimizer.step()}), the verifier catches $8/8$. On a
$50$-model false-positive sweep on clean models the verifier
reports $0$ spurious refutations.

\subsection{Localisation tracer and hybrid mode}
\label{sec:impl-loc-hybrid}

\paragraph{Localisation.} A v4 review noted that of $6$
detections in early experiments, only $1$ landed within $\pm 5$
lines of the bug. The v4 verifier returns
\texttt{location.line=0} for bugs detected via the model-check
path (those that bypass per-step counterexamples). We add a
``first unsatisfied symbolic constraint'' tracer
(\texttt{src/v5/localization.py}): given a refutation, we
extract the offending tensor and dim variables from the bug
message and walk the source AST to the earliest assignment or
typing-rule call that bound a symbolic dim appearing in the
unsatisfied constraint, attributing the bug to its source line.

\paragraph{Hybrid mode (TG-static $\to$ FakeTensor fallback).}
\texttt{tensorguard.hybrid\_check(M)}
(\texttt{src/v5/hybrid\_mode.py}) returns the TG verdict if it
is \textsc{Verified} or \textsc{Refuted}; otherwise, if $M$ is
importable and instantiable with a no-arg constructor, it falls
back to \texttt{FakeTensorMode}. Both the localisation tracer
and the hybrid mode are evaluated in
\Cref{sec:eval-loc-hybrid}.
